A revisão da literatura foi conduzida para contextualizar as escolhas de modelagem
adotadas neste trabalho e verificar, nos trabalhos encontrados, como o problema de
gestão de portfólio de contratos de energia tem sido abordado. O levantamento
permite identificar quais aspectos já foram tratados, quais simplificações foram
adotadas e em que pontos a formulação proposta neste trabalho se diferencia dos
estudos anteriores identificados.

A revisão apresentada neste capítulo foi conduzida a partir de buscas nas bases Portal de Periódicos CAPES, Google Scholar e repositórios institucionais de universidades brasileiras, utilizando as palavras-chave \emph{otimização de portfólio de energia}, \emph{comercialização de energia elétrica}, \emph{programação estocástica multiestágio}, \emph{stochastic dual dynamic programming} e \emph{gestão de contratos bilaterais}. Os trabalhos selecionados foram agrupados em três eixos temáticos: (i) otimização de portfólio de contratos de energia, (ii) métodos de solução para programação estocástica multiestágio aplicados a energia, e (iii) modelagem estocástica do PLD. A seção final posiciona o presente trabalho em relação às lacunas identificadas.

\section{Otimização de Portfólio de Contratos de Energia}
\label{sec:rl_portfolio}

A literatura sobre gestão de portfólio no setor elétrico brasileiro é predominantemente inspirada na teoria de Markowitz, transposta do mercado financeiro para o contexto de contratos de energia. Nessa abordagem, o problema consiste em determinar a alocação ótima entre diferentes tipos de contratos, incluindo bilaterais de prazo fixo, contratos no mercado \emph{spot} e, em alguns casos, contratos com fontes renováveis, de modo a maximizar o retorno esperado para um dado nível de risco.

\citeonline{gunn2012} aplica o modelo de Markowitz à comercialização de energia proveniente de novos empreendimentos, como biomassa, eólica e pequenas centrais hidrelétricas, no mercado brasileiro, modelando as incertezas de oferta e as regras de mercado por meio de simulações que alimentam a fronteira eficiente risco-retorno. \citeonline{guerreiro2017} avança nessa direção ao questionar a adequação da variância como medida de risco para o setor elétrico: dado que a distribuição do PLD é assimétrica e de cauda pesada, penalizar simetricamente ganhos e perdas extremas distorce a decisão de contratação. O autor propõe o CVaR como medida mais coerente, aplicando-o à otimização de portfólios de comercialização no ACL brasileiro. \cite{munhoz2018} segue linha semelhante, analisando o nível ótimo de contratação a preço fixo versus exposição ao PLD para um consumidor livre, com minimização de custo e risco via Markowitz.

Uma limitação comum a esses três trabalhos é a formulação estática do problema: o portfólio é otimizado em um único estágio, com base em uma distribuição de cenários gerada \emph{a priori}, sem modelagem explícita da sequência de decisões ao longo do tempo. Essa simplificação ignora o fato de que contratos bilaterais têm duração superior a um mês, comprometendo recursos financeiros e posições físicas em estágios futuros, e que novas informações sobre o PLD são reveladas progressivamente, permitindo ao agente adaptar sua estratégia de contratação.

\citeonline{josemaria2018} e \cite{carvalho2021} abordam o problema sob uma perspectiva mais descritiva, detalhando os tipos de contratos disponíveis no ACL, os riscos associados a cada modalidade e as estratégias adotadas por agentes na prática. \cite{carvalho2021} destaca que a literatura acadêmica tende a focar em modelos de previsão de preços, mas falha ao não incorporar explicitamente o perfil de aversão ao risco dos agentes reais, observação que reforça a necessidade de modelos que tratem a incerteza de forma estruturada.

O trabalho de \citeonline{pedrini2019} representa o avanço mais próximo da abordagem adotada neste trabalho dentro da literatura nacional: a autora formula a decisão de contratação de um consumidor livre como um problema estocástico de dois estágios, com a possibilidade de contratar geração eólica, e resolve o modelo por decomposição de Benders. A formulação de dois estágios captura a distinção entre decisões tomadas antes e após a revelação das incertezas, incorporando a não antecipatividade de forma estrutural. No entanto, o horizonte de dois estágios é uma simplificação para o problema de uma comercializadora, cujos contratos têm duração de vários meses e cujas decisões de contratação se repetem mensalmente ao longo de todo o ano.

Em síntese, a literatura nacional sobre portfólio de contratos de energia avançou de formulações estáticas baseadas em média-variância para abordagens com medidas de risco mais adequadas à assimetria do PLD e, num caso, para a estrutura de dois estágios. No entanto, formulações que combinem decisão sequencial ao longo de múltiplos meses com não antecipatividade explícita e restrições financeiras ao longo de todo o horizonte, no contexto da comercializadora no ACL brasileiro, ainda carecem de representação na literatura.


\section{Métodos de Solução para Programação Estocástica Multiestágio}
\label{sec:rl_sddp}

A extensão natural de dois para múltiplos estágios introduz um desafio computacional central: a abordagem mais direta de solução, a Formulação Determinística Equivalente (DEQ), que representa explicitamente toda a árvore de cenários em um único programa linear, tem porte que cresce exponencialmente com o número de estágios e ramos, conforme discutido na Seção~\ref{sec:prog_estocastica}. Para horizontes de doze meses com múltiplos cenários por período, a DEQ rapidamente excede os limites de memória disponíveis, tornando-se computacionalmente inviável.

O método de Programação Dinâmica Dual Estocástica (PDDE), introduzido por \citeonline{pereira1991} para o planejamento de operação de sistemas hidrotérmicos, contorna esse problema ao substituir a função de custo futuro por uma aproximação linear por partes construída iterativamente a partir de cortes de Benders. O tamanho de cada subproblema de estágio permanece fixo e independente do horizonte total, tornando o método escalável para instâncias de grande porte.

\cite{shapiro2011} analisa formalmente as propriedades estatísticas e a convergência do PDDE, estabelecendo as condições sob as quais os cortes acumulados aproximam a função de valor ótima. \cite{shapiro2013} estendem o método para formulações com aversão a risco, incorporando medidas como CVaR na função objetivo, e apresentam estudos computacionais relacionados ao planejamento do sistema interligado brasileiro, demonstrando a aplicabilidade do método em escala real. \cite{homemdemello2011} contribuem com uma análise prática sobre estratégias de amostragem e critérios de parada para o PDDE, avaliando o compromisso entre qualidade da solução e custo computacional no contexto do despacho hidrotérmico de longo prazo.

\cite{philpott2012} incorporam medidas de risco coerentes ao PDDE e comparam políticas com diferentes níveis de aversão ao risco em um sistema hidrotérmico, mostrando que a escolha da medida de risco afeta significativamente as decisões de despacho e o perfil de exposição do agente. \cite{street2020} tratam de uma questão de modelagem relevante para a estrutura de subproblemas do PDDE: a ordem entre a tomada de decisão e a revelação da incerteza dentro de cada período. Os autores mostram que a simplificação \emph{Hazard-Decision}, em que as decisões são tomadas assumindo conhecimento perfeito das afluências do período corrente, introduz distorções econômicas mensuráveis no despacho hidrotérmico brasileiro, reforçando a importância de modelar explicitamente a precedência entre decisão e realização da incerteza.

Do ponto de vista computacional, \citeonline{dowson2021} apresentam o \texttt{SDDP.jl}, pacote em Julia que implementa o PDDE sobre o ambiente de modelagem JuMP. O pacote organiza o problema como um grafo de política, uma generalização da árvore de cenários que acomoda estruturas mais flexíveis, incluindo a decomposição de cada período em nós de decisão e nós de liquidação, e gerencia a geração de cortes, os passes \emph{forward} e \emph{backward} e o cálculo dos limites de convergência. É a ferramenta utilizada na implementação da PDDE neste trabalho.

Em síntese, a PDDE é uma técnica consolidada para problemas de energia com estrutura multiestágio, com aplicações que vão do despacho hidrotérmico ao gerenciamento de risco com medidas coerentes. No entanto, sua aplicação à decisão sequencial de compra e venda de contratos bilaterais por uma comercializadora, com comparação sistemática frente ao DEQ em diferentes portes de instância, permanece pouco explorada na literatura.


\section{Modelagem Estocástica do PLD}
\label{sec:rl_pld}

A literatura identifica um conjunto de propriedades do PLD que tornam inadequada sua representação por um único valor esperado: sazonalidade marcada pelo ciclo hidrológico, dependência temporal entre períodos consecutivos, distribuição assimétrica com caudas pesadas e limites regulatórios de piso e teto \cite{guerreiro2017,ccee2024a,mme2002}. \citeonline{guerreiro2017} documenta que a assimetria positiva da distribuição, com picos pronunciados nos períodos secos, leva modelos baseados em valor esperado a subestimar sistematicamente o custo de posições deficitárias, reforçando a necessidade de modelos estocásticos com sazonalidade representada de forma estrutural.

\citeonline{reis2013} fornece a referência metodológica central para essa classe de modelos: os Modelos Autorregressivos Periódicos PAR($p$), que atribuem um conjunto distinto de parâmetros autorregressivos para cada mês do ano, acomodando a estacionariedade periódica das séries hidrológicas e de energia. O modelo GEVAZP, desenvolvido pelo CEPEL para geração de séries sintéticas de energia natural afluente no planejamento do SIN, baseia-se nessa mesma classe. \cite{detzel2011} validam um modelo AR(1) parcimonioso para afluências em usinas da região sul do Brasil, confirmando que modelos autorregressivos de baixa ordem capturam a estrutura de dependência temporal relevante sem requerer estimação de muitos parâmetros.

A utilidade desses modelos autorregressivos está, precisamente, em sua capacidade de gerar cenários que alimentem os métodos de otimização. Para o DEQ, os cenários formam a árvore explícita de realizações sobre a qual o problema é resolvido \cite{birge2011}. Para a PDDE, os cenários são amostrados nas trajetórias da etapa \emph{forward} do algoritmo, sem construção explícita da árvore completa \cite{pereira1991,shapiro2009}. Em ambos os casos, a adequação estatística dos cenários à distribuição real do PLD condiciona diretamente a qualidade da política obtida.

Considerando esses requisitos de adequação estatística, o arcabouço PAR(1) é compatível com o problema de portfólio tratado neste trabalho: captura a sazonalidade mensal e a dependência temporal com apenas doze coeficientes autorregressivos, pode ser estimado a partir de séries históricas da CCEE, e alinha-se com a tradição do GEVAZP validada no contexto brasileiro \cite{reis2013,detzel2011}. A estimação em escala logarítmica garante positividade das séries sintéticas, compatível com a natureza do PLD como preço com piso regulatório positivo.


\section{Posicionamento do Trabalho e Lacunas Identificadas}
\label{sec:rl_posicionamento}

A revisão da literatura permite identificar três lacunas que motivam e justificam o presente trabalho.

A primeira diz respeito à profundidade da formulação temporal. Os trabalhos de \cite{guerreiro2017}, \cite{gunn2012} e \cite{munhoz2018} tratam o problema de forma estática, sem modelagem da interdependência entre decisões de contratação de diferentes meses. \citeonline{pedrini2019} avança para dois estágios, mas o horizonte é insuficiente para capturar a dinâmica de contratos com duração de vários meses. Na literatura revisada, não foram identificados trabalhos que combinem, em um único modelo, a formulação multiestágio de decisões mensais de contratação, a não antecipatividade explícita e restrições de fluxo de caixa ao longo de todo o horizonte de planejamento para o contexto específico da comercializadora no ACL brasileiro.

A segunda lacuna diz respeito à comparação empírica entre abordagens de solução. A superioridade teórica da decomposição sobre a formulação explícita em termos de escalabilidade é bem estabelecida desde \cite{pereira1991}. Contudo, na literatura revisada sobre gestão de portfólio de contratos bilaterais no mercado brasileiro, não foram identificados trabalhos que documentem essa comparação empiricamente --- avaliando os limites concretos de cada abordagem em termos de memória consumida, tempo de processamento e qualidade da solução para instâncias de diferentes portes.

A terceira lacuna diz respeito à modelagem da solvência financeira intra-horizonte. Os trabalhos revisados focam predominantemente no balanço de energia --- posição física líquida e exposição ao PLD --- sem incluir restrições explícitas de fluxo de caixa e limite de crédito que condicionem as decisões de contratação ao longo do horizonte. Para uma comercializadora sem geração própria, essa dimensão financeira é tão relevante quanto o balanço físico, pois contratos de compra exigem desembolsos futuros que podem comprometer a liquidez da empresa.

O presente trabalho endereça as três lacunas ao formular o problema de gestão de portfólio da comercializadora como um programa estocástico multiestágio com restrições explícitas de fluxo de caixa, resolvido tanto por DEQ quanto por PDDE, com comparação empírica sistemática entre as duas abordagens em instâncias sintéticas parametrizadas a partir de dados históricos do SIN.

